% arXiv source: standard packages and vector PDF figures; no shell escape.
% House style follows output/pdf/conv_paper_fused.tex (IEEEtran journal).
\documentclass[journal]{IEEEtran}
\usepackage[T1]{fontenc}
\usepackage{lmodern}
\usepackage{amsmath,amssymb,amsthm}
\usepackage{bm}
\usepackage{microtype}
\usepackage{graphicx}
\usepackage{booktabs}
\usepackage[hidelinks]{hyperref}
\setlength{\textfloatsep}{8pt plus 2pt minus 2pt}
\setlength{\floatsep}{7pt plus 2pt minus 2pt}
\setlength{\intextsep}{7pt plus 2pt minus 2pt}
\newtheorem{definition}{Definition}
\newtheorem{theorem}{Theorem}
\newtheorem{proposition}{Proposition}
\newtheorem{corollary}{Corollary}
\newtheorem{lemma}{Lemma}
\newcommand{\R}{\mathbb{R}}
\newcommand{\E}{\mathbb{E}}
\newcommand{\Cov}{\operatorname{Cov}}
\newcommand{\diag}{\operatorname{diag}}
\newcommand{\tr}{\operatorname{tr}}
\newcommand{\norm}[1]{\left\lVert#1\right\rVert}
\newcommand{\trans}{\mathsf{T}}
\newcommand{\given}{\,\middle|\,}
\newcommand{\Normal}{\mathcal{N}}
\newcommand{\data}{\mathcal{D}}
\newcommand{\clip}{\operatorname{clip}}
\makeatletter
\newcommand{\inputrows}[1]{\@@input #1 }
\makeatother
\graphicspath{{figures/}}
\title{Exact Markov Elimination for Relational Current Tracking}
\author{Anonymous}
\date{}
\hypersetup{pdftitle={Exact Markov Elimination for Relational Current Tracking},pdfauthor={Anonymous}}
\begin{document}
\maketitle

\begin{abstract}
Glory first be to The Holy One, Blessed be He, from whom all good things
come. We study position-only tracking with a finite Gaussian current prior
and an observation-driven acquisition rate. Position is the exact integral
of orthonormal cell currents; a five-component correlation-length mixture
retains uncertainty about their continuation. We derive a four-state
realization per spatial coordinate and mixture component that reproduces
the original finite midpoint covariance, noisy anchor, observation
likelihood, and posterior. This eliminates dense current-history covariance
from streaming updates without changing the represented trajectory model.
On an M4 CPU, a 128-cell observation update decreases from 452.7 to
12.2 microseconds; the corresponding NumPy realization already achieves
23.4 microseconds. Across 480 matched acquisition cases, the implementations
produce identical acquisition times and boundary-trigger decisions, with
position and forecast differences below $9\times10^{-12}$ meters. A
separate retained synthetic comparison includes constant-velocity and
constant-acceleration Kalman filters, interacting multiple-model filters,
a robust variant, and a three-dimensional turn unscented Kalman filter.
At low noise with 65 observations, the current model's two-second endpoint
RMSE is 2.958 meters, versus 3.269 for robust IMM; the advantage does not
persist decisively in every noise and sampling condition. Adaptive density
reduces acquisition count and improves tracking after a settling interval.
We also give explicit model-conditional uncertainty growth and a sharp
bounded-acceleration error certificate common to all compared estimators;
no superior worst-case breakdown rate is established.
Exact elimination changes the computational balance enough that scheduling,
rather than assimilation, becomes the dominant adaptive-loop cost.
\end{abstract}

\begin{IEEEkeywords}
Target tracking, Gaussian processes, Kalman filtering, Mat\'ern covariance,
state-space realization, adaptive sampling, exact elimination.
\end{IEEEkeywords}

\section{Introduction}

A position tracker uses accumulated observations to correct its current
belief and to continue that belief beyond the last measurement. These two
operations depend on the same temporal law. A representation that exposes
all past and future current coefficients makes that dependence explicit,
but it can conceal a much smaller sufficient filtering state. The
computational question is therefore precise: can the calculation be
reorganized while preserving the prior, evidence, and resulting decisions?

We consider a finite current construction on a declared time interval.
Each Cartesian position is an anchor plus the exact integral of a
piecewise-constant current. The current covariance combines a shared
constant component with a Mat\'ern-$3/2$ component, and a fixed mixture of
five correlation lengths represents uncertainty about continuation scale.
The first implementation conditioned a dense covariance in either real
current coordinates or a unitary Zak representation. Both coordinates
produce the same posterior when the full covariance is retained.

The present work derives and implements a four-state realization of this
\emph{finite} prior. Two states realize the sampled Mat\'ern current; the
other two carry the shared constant current and accumulated position.
The resulting filter retains every length hypothesis and its exact
Gaussian evidence. The reformulation changes the order in which latent
history is eliminated. It does not replace the finite current by a
continuously integrated process, discard mixture components, or introduce
a new correction-selection rule.

There are three distinct empirical questions. First, how does the declared
current prior compare with existing tracking implementations? Second,
what accuracy and acquisition-count tradeoff follows when a boundary
estimate increases observation density? Third, what computational cost
can be removed from exactly the same estimation calculation? We answer
these questions with separate protocols and retained records. We do not
combine their timings or error statistics into a single ranking.

The contributions are a finite-model equivalence theorem including a
noisy anchor and a severed-current ablation; an implementation that shares
spatial covariance, evidence normalization, and cached cell propagation;
a matched evaluation of numerical equivalence and acquisition cost; and
explicit uncertainty-growth bounds that separate model-conditional
posterior statements from physical prediction-error certificates.
The empirical comparison supports competitiveness on the tested
synthetic families, with uncertainty intervals and failure-to-dominate
conditions reported alongside the favorable case. No real-flight or
universal tracking-optimality claim is made.

\section{Relation to Existing Tracking Methods}

Kalman filtering supplies recursive inference for linear Gaussian state
models~\cite{kalman1960}. A temporal Gaussian process can also admit a
finite Markov realization. Hartikainen and S\"arkk\"a give exact
state-space constructions for temporal Mat\'ern covariance functions and
discuss nonuniform observations~\cite{hartikainen2010}. Parallel temporal
GP algorithms further exploit associative filtering and smoothing
structure~\cite{corenflos2021}. These are established connections. The
four-state realization below specializes them to a fixed midpoint-current
model and its cellwise integral; we do not claim the general
GP--state-space equivalence as new.

An interacting multiple-model filter (IMM) instead mixes dynamically
switching motion hypotheses~\cite{blom1988}. Its model set and transition
law determine which maneuvers can be represented. Unscented filtering
propagates nonlinear state uncertainty through sigma points
\cite{julier2004}. Our measured turn UKF estimates a continuous
three-dimensional turn vector, whereas our IMM uses a finite catalog of
axis-aligned turn generators in its state coordinates. This makes the
UKF a useful control for the IMM catalog's limitations.

Current tracking research includes comparisons of fixed and adaptive
multiple-model transition laws~\cite{lizzio2026}, IMM filtering and
smoothing on manifolds~\cite{koller2021}, and feedback particle filters on
matrix Lie groups~\cite{zhang2017}. Their observation contracts and
algorithms differ from the local comparators evaluated here. In
particular, our robust IMM is a specified contamination-likelihood
implementation, not a reproduction of every adaptive robust IMM method.
The recent studies provide context, not additional measured baselines.

The length mixture in this paper has a \emph{fixed latent length per
trajectory hypothesis}. Its weights update by marginal likelihood; its
states are not mixed through an IMM transition matrix. It consequently
represents uncertainty about temporal correlation rather than a Markov
sequence of switching maneuver labels. ``Relational'' refers here to the
specified covariance between cell currents. No general nonlinear
relational denoiser is implied by that name.

\section{Finite Current and Observation Model}
\label{sec:model}

Let $[0,T]$ be a declared representation interval, $N$ the number of equal
cells, and $h=T/N$. The endpoint $T$ is known to the representation; it is
not an observation of future motion. All queries lie in this interval.

\begin{definition}[Finite current position]
For one spatial coordinate, let $a$ be an anchor error and $v_i$ the
current value in cell $i$, $0\leq i<N$. Relative to the first measured
position $y_0$, define
\begin{equation}
 x(t)=y_0+a+\sum_{i=0}^{N-1}v_i\,\clip(t-ih,0,h).
 \label{eq:position}
\end{equation}
The orthonormal cell basis and its coefficients are
\begin{equation*}
 \phi_i(t)=h^{-1/2}\bm{1}_{[ih,(i+1)h)}(t),\qquad c_i=\sqrt h\,v_i.
\end{equation*}
The current is therefore $\sum_i c_i\phi_i$.
\end{definition}

For each correlation length $\ell$, let $\lambda_\ell=\sqrt{3}/\ell$ and
prescribe the stationary cell-current covariance
\begin{align}
 \Cov(v_i,v_j\mid\ell)
 &=q_g+q_w(1+\lambda_\ell d_{ij})e^{-\lambda_\ell d_{ij}},
 \label{eq:kernel}\\
 d_{ij}&=|i-j|h.\nonumber
\end{align}
All experiments fix $q_g=q_w=4\,\mathrm{m}^2\mathrm{s}^{-2}$ and
\begin{equation}
 \ell\in\mathcal L=\{0.3,0.7,1.5,3,6\}\ \mathrm{s},
 \qquad \Pr(\ell)=1/5.
 \label{eq:lengths}
\end{equation}
The covariance is evaluated at cell midpoints. Their common half-cell
offset cancels in $d_{ij}$. The coefficient covariance is $h$ times
\eqref{eq:kernel}; increasing $N$ does not silently rescale current units.

Observations are complete three-dimensional positions
\begin{equation}
 \bm y_j=\bm x(t_j)+\bm\epsilon_j,
 \qquad \bm\epsilon_j\sim\Normal(0,\sigma^2I_3),
 \label{eq:observations}
\end{equation}
with known $\sigma>0$ and independent errors. The history begins at
$t_0=0$ and has strictly increasing arrival times. Conditioning a flat
absolute-position prior on $\bm y_0$ gives
$\bm a\sim\Normal(0,\sigma^2I_3)$ relative to that observation. This
anchor remains uncertain. The later sensor errors are independent of it;
$\bm y_0$ is not assimilated a second time. Conditional on $\ell$, the
three axes have independent, identical temporal priors.

A dense coefficient state has $N+1$ entries per axis. Its measurement row is
\begin{equation}
 B(t)=\left[1,\left\{\frac{\clip(t-ih,0,h)}{\sqrt h}\right\}_{i=0}^{N-1}\right].
 \label{eq:dense-row}
\end{equation}
An invertible coordinate change preserves the observation law when the
state, covariance, and measurement rows all transform consistently. In
particular, a unitary Zak matrix $U$ gives covariance $UPU^*$ and row
$BU^*$. The complex representation must retain the real-signal conjugacy
constraints; it does not justify a proper-complex observation likelihood.

\section{Exact Four-State Realization}
\label{sec:realization}

Write $v_i=g+w_i$, with independent
$g\sim\Normal(0,q_g)$ and a Mat\'ern-$3/2$ current sequence $w_i$.
Introduce the normalized derivative $r=w'/\lambda_\ell$ of an auxiliary
continuous process, sampled at the current midpoints. This auxiliary
process realizes the covariance; the position remains the finite
integral in \eqref{eq:position}.

\begin{lemma}[Sampled current realization]
\label{lem:current}
Put $\alpha=\lambda_\ell h$ and
\begin{equation}
 A_\ell=e^{-\alpha}
 \begin{bmatrix}1+\alpha&\alpha\\-\alpha&1-\alpha\end{bmatrix},
 \quad Q_\ell=q_w(I_2-A_\ell A_\ell^\trans).
 \label{eq:A}
\end{equation}
For $z_i=(w_i,r_i)^\trans$, the stationary recursion
\begin{equation}
 z_{i+1}=A_\ell z_i+\eta_i,\qquad
 \eta_i\sim\Normal(0,Q_\ell),\qquad
 z_0\sim\Normal(0,q_wI_2)
 \label{eq:z}
\end{equation}
has the Mat\'ern part of \eqref{eq:kernel} as its first-coordinate
covariance. The innovations are independent across cells.
\end{lemma}
\begin{proof}
The continuous generator
$D=\lambda_\ell\left[\begin{smallmatrix}0&1\\-1&-2\end{smallmatrix}\right]$
satisfies
\begin{equation}
 D(q_wI_2)+(q_wI_2)D^\trans
 =-\diag(0,4\lambda_\ell q_w).
\end{equation}
Thus $Q_\ell$ is the integral of
$e^{Du}\diag(0,4\lambda_\ell q_w)e^{D^\trans u}$ over $0\leq u\leq h$
and is positive semidefinite. Exponentiating $Dh$ gives \eqref{eq:A}.
The nilpotent part of $D+\lambda_\ell I_2$ squares to zero, so
\begin{equation}
 A_\ell^n=e^{-n\alpha}
 \begin{bmatrix}1+n\alpha&n\alpha\\-n\alpha&1-n\alpha\end{bmatrix}.
 \label{eq:An}
\end{equation}
Stationarity follows from $A_\ell(q_wI_2)A_\ell^\trans+Q_\ell=q_wI_2$.
The $(1,1)$ entry of $q_wA_\ell^n$ is
$q_w(1+n\alpha)e^{-n\alpha}$, as required.
\end{proof}

Let $p_i$ be the accumulated position relative to $y_0$ at the start of
cell $i$. Define $s_i=(p_i,g,w_i,r_i)^\trans$. One cell advances by
\begin{align}
 s_{i+1}&=F_\ell s_i+\xi_i,\label{eq:state}\\
 F_\ell&=\begin{bmatrix}
 1&h&h&0\\0&1&0&0\\
 0&0&(A_\ell)_{11}&(A_\ell)_{12}\\
 0&0&(A_\ell)_{21}&(A_\ell)_{22}
 \end{bmatrix},\nonumber\\
 \Cov(\xi_i)&=\diag(0,0,Q_\ell),\nonumber\\
 \Cov(s_0)&=\diag(\sigma^2,q_g,q_w,q_w).
 \label{eq:statecov}
\end{align}
Here the block notation in \eqref{eq:statecov} places $Q_\ell$ in the
last two coordinates. At $t\in[ih,(i+1)h]$, let $\tau=t-ih$. The observation
row and position are
\begin{equation}
 H(t)=[1,\tau,\tau,0],\qquad x(t)=y_0+H(t)s_i.
 \label{eq:H}
\end{equation}
At $t=T$, use $i=N-1$ and $\tau=h$. Multiple observations within a cell
use the same latent state with different rows; no artificial current
innovation is introduced between them.

\begin{theorem}[Finite-model equivalence]
\label{thm:equivalence}
For fixed $N,T,\sigma$ and $\mathcal L$, the dense model
\eqref{eq:position}--\eqref{eq:kernel} and the realization
\eqref{eq:state}--\eqref{eq:H} define identical joint position laws conditional on the initial
observation $\bm y_0$.
After any finite increasing observation history satisfying
\eqref{eq:observations}, they therefore have identical component
posterior laws, marginal likelihoods, mixture weights, and position
posterior laws at all represented query times.
\end{theorem}
\begin{proof}
Lemma~\ref{lem:current}, together with the independent $g$, gives exactly
\eqref{eq:kernel}. The initialization gives the same independent anchor
error. The first row of $F_\ell$ yields
$p_i=a+h\sum_{j<i}(g+w_j)$, so \eqref{eq:H} is exactly
\eqref{eq:position}. These are equal linear functions of equal Gaussian
prior laws. Adding the same independent observation errors preserves
the joint law of observations and queries. Gaussian conditioning gives
the same conditional distributions and evidence for each $\ell$.
Bayes' rule with the same length prior then gives the same mixture.
\end{proof}

\begin{corollary}[Filtering sufficiency]
Given the component posterior over $s_i$ and the length weights after the
latest observation, all current and future position marginals can be
computed without retaining the dense coefficient posterior.
\end{corollary}
\begin{proof}
Future current innovations in \eqref{eq:z} are independent of the past
conditional on $z_i$. The state also contains the accumulated position
and shared constant current. Conditioning on these variables separates
future propagation from eliminated history.
\end{proof}

This theorem concerns the finite model at a fixed $N$. It is not a claim
that changing $N$ leaves the prior unchanged, nor that a continuous
Mat\'ern integral has the same covariance. The auxiliary derivative in
Lemma~\ref{lem:current} supplies the Markov realization while the finite
cell integral supplies the actual position law.

\subsection{Severed Current Covariance}

For an ablation time $t_s$, the dense reference zeros covariance between
currents whose midpoints lie on opposite sides of $t_s$. Let
$k_s=\#\{i:(i+1/2)h<t_s\}$. After integrating cell $k_s-1$, reset
$(g,w,r)$ to an independent Gaussian with covariance
$\diag(q_g,q_w,q_w)$, preserving accumulated position. The reset replaces
cross-group current covariance by zero and reproduces the dense ablation
exactly. For $k_s=0$ or $N$, there is no internal reset. This operation
removes a continuation connection deliberately; it is not an alternative
estimate selected from future truth.

\section{Shared Conditioning and Compiled Evaluation}
\label{sec:implementation}

Let $M_\ell\in\R^{4\times3}$ contain the component state means and
$P_\ell\in\R^{4\times4}$ its common axis covariance immediately before
an observation. Compute
\begin{align}
 b_\ell&=P_\ell H^\trans,&S_\ell&=Hb_\ell+\sigma^2,\nonumber\\
 e_\ell&=\bm y-\bm y_0-HM_\ell,\label{eq:innovation}\\
 M_\ell^+&=M_\ell+(b_\ell/S_\ell)e_\ell,&
 P_\ell^+&=P_\ell-b_\ell b_\ell^\trans/S_\ell.
 \label{eq:update}
\end{align}
The innovation $e_\ell$ is a row vector. The same rank-one covariance
update serves all three axes. The log-evidence increment is
\begin{equation}
 \Delta L_\ell=-\tfrac12
 \left(3\log(2\pi S_\ell)+\norm{e_\ell}^2/S_\ell\right).
 \label{eq:evidence}
\end{equation}
Initialize $L_\ell=-\log|\mathcal L|$ and normalize once after each
observation by $\pi_\ell=e^{L_\ell}/\sum_k e^{L_k}$, using log-sum-exp.
Readouts reuse the weights and total log evidence.

For component query means $\mu_\ell(t)$ and scalar variances $V_\ell(t)$,
the reported moments are
\begin{align}
 \bar\mu(t)&=\sum_\ell\pi_\ell\mu_\ell(t),\label{eq:mixmean}\\
 \Sigma(t)&=\sum_\ell\pi_\ell
 \left[V_\ell(t)I_3+\delta_\ell(t)\delta_\ell(t)^\trans\right],
 \quad \delta_\ell=\mu_\ell-\bar\mu.
 \label{eq:mixcov}
\end{align}
The between-component term is retained. The result need not be spatially
isotropic even though each component's within-axis covariance is shared.

\begin{proposition}[Likelihood refinement and coordinates]
Assimilating one Gaussian likelihood in $m$ equal fractional powers gives
the same posterior as assimilating it once. An invertible state-coordinate
change with consistent covariance and measurement transforms gives the
same observation evidence and physical position posterior.
\end{proposition}
\begin{proof}
For a Gaussian likelihood $\mathcal{L}$, the product of $m$ factors
$\mathcal{L}^{1/m}$ equals $\mathcal{L}$. Each normalized Gaussian factor
has noise variance $m\sigma^2$; its normalization differs from that of
$\mathcal{L}^{1/m}$ only by a state-independent constant. Posterior
normalization removes that constant. Model evidence must still be
computed under the original likelihood, as in \eqref{eq:evidence},
rather than by treating the factors as additional physical observations.
A coordinate change leaves the physical observation random variable
unchanged, proving the second statement.
\end{proof}

Integer-cell propagation uses shared precomputed tables
\begin{align}
 F_{\ell,n}&=F_\ell^n,&Q_{\ell,0}&=0,\nonumber\\
 Q_{\ell,n+1}&=F_\ell Q_{\ell,n}F_\ell^\trans+Cov(\xi_i).
 \label{eq:jump}
\end{align}
A jump of $n$ cells applies $F_{\ell,n}$ and $Q_{\ell,n}$ once. Each
query propagates a copy of the filtering state, so evaluation does not
change subsequent estimates or acquisition decisions. The compiled
kernel fuses branch propagation, rank-one updates, and evidence
normalization into a single update call; a separate call evaluates all
requested current/future marginal moments. Double precision is used
without fast-math.

With $K=|\mathcal L|$ fixed, streaming update work is $O(K)$ after the
$O(KN)$ table construction. At $N=128$, covariance entries per branch
decrease from $129^2=16641$ to $4^2=16$. The five filtering means and
covariances occupy 1120 bytes in double precision, excluding weights,
object overhead, and shared tables. This is a statement about the
filtering state, not total program memory.

The implementation retains the observation history for compatibility.
Queries earlier than the last observation, and exports of the full
coefficient posterior, lazily replay the dense reference. Their cache
is invalidated by a new observation. These operations retain dense costs.
The transition cache is bounded and shared between trackers, and the
retained input history grows with the observation count. The fast API
returns position marginals; it does not return a certified continuous
boundary-crossing probability from those marginals alone.

\section{Causal Acquisition Density}
\label{sec:acquisition}

The acquisition experiment changes sampling cadence while retaining one
posterior and one update law. At each acquired position, the tracker is
updated once. The first acquisition time for which the updated estimated
first coordinate is nonnegative becomes a latched trigger $t_\star$.
No true crossing time enters this trigger. For $t\geq t_\star$, the
requested rate is
\begin{equation}
 r(t)=r_0+(r_1-r_0)\left(1-e^{-(t-t_\star)/\tau_r}\right),
 \label{eq:rate}
\end{equation}
with $r_0=2$ Hz, $\tau_r=1$ s, and
$r_1\in\{4,8,16,32\}$ Hz. Before the trigger $r(t)=r_0$.
Two controls use constant 2 and 32 Hz rates.

Irregularity is introduced in phase rather than by rounding times to a
common grid. Draw $u_j\sim\mathrm{Uniform}(0.8,1.2)$ independently and
choose the next time from
\begin{equation}
 \int_{t_j}^{t_{j+1}}r(s)\,ds=u_j.
 \label{eq:phase}
\end{equation}
The retained implementation uses 44 bisection steps for the ramped
integral. This scheduler is unchanged in the optimized replay. The rate
command is continuous, but individual observations and posterior updates
remain discrete. The covariance-weighted innovation in
\eqref{eq:update} supplies the correction; no separate witness gate or
boundary-dependent filter selection is applied.

\begin{proposition}[Causal policy equivalence]
For the same initial observation and random phase/noise streams, exact
posterior equivalence implies identical acquisition sequences under
\eqref{eq:rate}--\eqref{eq:phase}, provided the same numerical scheduling
map is used in the exact-arithmetic comparison.
\end{proposition}
\begin{proof}
The policies initially use the same observation and schedule. If the
acquired history agrees through step $j$, Theorem~\ref{thm:equivalence}
gives the same corrected position. The trigger condition and rate
therefore agree. The same next phase gives the same next time, and the
same observation law and noise draw give the same next observation.
Induction proves the claim.
\end{proof}

Floating-point reordering can change a sign decision exactly at a
threshold. Consequently exact decision agreement is measured as well as
proved algebraically; no universal bitwise claim follows from the theorem.

\section{Experimental Protocols}
\label{sec:protocols}

\subsection{Matched Forecast Comparison}
\label{sec:forecast-protocol}

The retained screen contains five motion families: straight motion,
tilted helices, spatial figure eights, smooth stop/restart, and motion
approaching changing waypoints. Truth uses analytic positions or a
separate fine-step Cartesian simulation. Random orientations and
translations avoid a preferred global observation frame. Future
waypoint commands are unavailable to all estimators.

Six fresh seeds per family give 30 trajectories in each of three paired
conditions: C1 uses 65 observations with $\sigma=0.35$ m; C2 uses 65 with
$\sigma=1.4$ m; C3 uses 9 with $\sigma=0.35$ m. Histories span $[0,6]$ s;
queries are 20 times from 6.1 through 8 s. The 65 observations are evenly
spaced; the sparse condition takes every eighth history observation.
Noise draws are paired across conditions. The current representation
uses $T=8$ s and $N=64$. Eight methods give 720 method/case runs.
Future truth is used only by the evaluator.

The five conventional controls are local literature-family
implementations. CV and CA use linear kinematics; IMM mixes CV, CA,
and six fixed turns of magnitude $0.65$ rad/s about signed Cartesian
axes. Its mixing time scale is 3 s. Robust IMM uses the same model set
with a two-component observation-noise mixture having probabilities
$(0.97,0.03)$ and variances $(\sigma^2,100\sigma^2)$, recombining the
noise branches by moments. The continuous-turn UKF uses
$(\alpha,\beta,\kappa)=(1,2,0)$ and estimates a full 3-D turn vector.
All conventional filters include the pre-existing augmented position/bias
state. Their priors are therefore not identical to the current prior.

Configurations were frozen from a separate ten-case development stage,
with three candidate configurations per method and objective tracking
MSE plus $0.25$ times two-second forecast MSE. Selected process
intensities are $q=1$ for CV and robust IMM, and $q=0.1$ for CA and IMM;
the UKF turn/speed diffusivities are $0.75/0.30$. Only supplied sensor
$\sigma$ changes in this screen. The five current lengths and current
variances were declared before inspecting this screen and were not
selected on its future truth. Unequal development histories and differing
model classes limit claims of globally optimal parameter matching.

The current model produces analytic forecast moments. Conventional
forecasts retain the existing 2048-path Monte Carlo implementation, with
separate fixed forecast seeds. Thus their estimated means include finite
Monte Carlo error. These historical forecast workloads are not used as
matched-kernel timing comparisons.

The primary error is the Euclidean 3-D position error at 8 s. Reported
RMSE is the square root of mean squared endpoint distance over the 30
cases. Paired percentile intervals use 4000 resamples of the six seeds
within each of the five fixed families. They are pointwise intervals,
not simultaneous tests or uncertainty over unseen motion families.

\subsection{Acquisition Accuracy and Matched Replay}

Boundary-crossing fixtures use the same five families with eight new
seeds and two noise levels. Their first coordinate is replaced by
$2(t-6)+0.3\sin(0.7t+\varphi)$, with seed-dependent phase; the remaining
coordinates follow the existing generated family. Truth is represented
on a fixed 0.01 s grid over $[0,14]$ s and linearly interpolated for
queries. These are explicitly modified crossing fixtures, not the
unchanged trajectories of Section~\ref{sec:forecast-protocol}.

The tracker uses $T=14$ s, $N=128$. Six acquisition policies give 480
policy/case runs. Current position and two-second forecasts are evaluated
at the same 0.1 s grid over $[0,12]$ s for every policy. The true crossing
is used only to define evaluation windows: the first second, seconds
3--6 after crossing, and the full first six post-crossing seconds,
intersected with the evaluation interval. Per-case MSEs are averaged
with equal case weight before taking their square root.

The exact-optimization study replays all 480 cases through the archived
dense implementation and compiled realization, giving 960 matched runs.
Run order alternates by family/seed to reduce systematic order effects.
Sample arrays, decisions, current estimates, forecasts, and uncertainty
radius diagnostics are compared. There is no timing-based or
accuracy-based tracker selection during acquisition.

\subsection{Cost Controls}

Microbenchmarks use an identical 160-observation irregular stream, 31
repeats after one warm-up, double precision, and one BLAS thread on the
M4 Mini CPU. Dense 512-cell configurations use five repeats. The four
implementations are dense Zak, dense real current, four-state NumPy,
and four-state compiled. Each repeat times updates, current-position
readouts, and two-query forecasts separately. Cold initialization clears
the transition cache; warm initialization is measured in the acquisition
replay. No CPU instruction counters were collected: the evidence is
algebraic state/work reduction and measured elapsed time.

A separate small CV cost control has two states per axis, a shared
isotropic covariance, white-acceleration intensity $q=1$, and no mixture
or bias state. It is implemented in NumPy and compiled C++, with the same
irregular observation stream and Python interface style. This control
establishes a simple implementation cost scale. It is not the augmented
CV model in the earlier accuracy screen and carries no matched-accuracy
claim.

Total online acquisition cost includes initialization, assimilation,
readouts, and scheduling. Diagnostic-only pre-update probes and offline
truth/error scoring are excluded. Sensor energy, communications,
bandwidth, and acquisition latency are not assigned a monetary or CPU
cost. All timing claims refer to these retained implementations and this
host rather than a real-time deadline guarantee.

\section{Uncertainty Growth and Conditional Error Bounds}
\label{sec:bounds}

Exact inference does not by itself bound error against an arbitrary
trajectory. We distinguish a posterior statement conditional on the
current prior, a physical error certificate conditional on motion and
sensor bounds, and empirical degradation under changed conditions.

\subsection{Growth Under the Declared Current Prior}

Fix a realized observation history ending at $t_b$ and a represented
query $t_b+u\leq T$. In component $\ell$, write the posterior position
variance at $t_b$ as $V_{b,\ell}$ per coordinate. Let $a_i(u)$ be the
length of overlap between cell $i$ and $[t_b,t_b+u]$, so
$0\leq a_i\leq h$ and $\sum_i a_i=u$. Define
\begin{align}
 W_\ell(u)&=q_g u^2+q_w\sum_{i,j}a_i a_j\rho_{|i-j|},\\
 \rho_k&=(1+k\alpha)e^{-k\alpha},\qquad\alpha=\lambda_\ell h.
\end{align}
Gaussian conditioning cannot increase the covariance of this fixed
linear increment. Cauchy--Schwarz on its posterior covariance with the
current position therefore gives
\begin{equation}
 V_\ell(t_b+u)\leq\bigl(\sqrt{V_{b,\ell}}+\sqrt{W_\ell(u)}\bigr)^2.
 \label{eq:variance-growth}
\end{equation}
The nonnegative Toeplitz kernel has row sum at most
\begin{equation}
 C_\alpha=1+2\left(\frac{z}{1-z}+\frac{\alpha z}{(1-z)^2}\right),
 \qquad z=e^{-\alpha}.
\end{equation}
Using $\rho_k\leq1$ and the row-sum operator bound with
$\sum_i a_i^2\leq hu$ proves
\begin{equation}
 W_\ell(u)\leq q_g u^2+q_w\min\{u^2,hC_\alpha u\}.
 \label{eq:increment-growth}
\end{equation}
Thus component standard deviation has an explicit at-most-linear
envelope in forecast horizon, within the declared interval. This is a
property of the chosen stochastic prior, not a bound on arbitrary motion.

Let $\mu_\ell(u)$ and $V_\ell(u)I_3$ be the component position moments,
and $\bar\mu(u)$ the mixture mean. For $0<\delta<1$, the ball of radius
\begin{equation}
 R_\delta(u)=\max_\ell\left\{\norm{\mu_\ell(u)-\bar\mu(u)}
 +\sqrt{6V_\ell(u)\log(6/\delta)}\right\}
 \label{eq:mixture-ball}
\end{equation}
about $\bar\mu(u)$ has posterior probability at least $1-\delta$
\emph{under the declared mixture model}. Indeed, the coordinate Gaussian
tail bound and a union bound give the stated component radius; the
triangle inequality places every component ball inside this ball.
Averaging the conditional probabilities over the length weights preserves
$1-\delta$. This includes between-component mean separation, unlike
using component variance alone. It is a single-query posterior statement,
not a simultaneous continuous-time tube or a guarantee for the synthetic
truth families. Exact elimination preserves this bound and its inputs.

\subsection{A Sharp Bound on Actual Prediction Error}

A separate, estimator-independent assumption supplies a physical
certificate. Suppose $\bm x'$ is absolutely continuous and
$\norm{\bm x''(t)}\leq A$ almost everywhere over the observation and
forecast interval. Suppose two observations at $t_a<t_b$ satisfy
$\norm{\bm y_i-\bm x(t_i)}\leq\varepsilon_i$. These assumptions do not
follow from a nominal sensor standard deviation or from the current prior.

\begin{theorem}[Two-observation continuation bound]
\label{thm:physical}
Put $\Delta=t_b-t_a>0$ and $u\geq0$. Define
\begin{align}
 \bm c(u)&=\bm y_b+\frac{u}{\Delta}(\bm y_b-\bm y_a),\\
 r(u)&=\left(1+\frac{u}{\Delta}\right)\varepsilon_b
       +\frac{u}{\Delta}\varepsilon_a+\frac A2u(\Delta+u).
 \label{eq:physical-radius}
\end{align}
Then $\norm{\bm x(t_b+u)-\bm c(u)}\leq r(u)$.
The constants are sharp even in one dimension.
For any estimator output $\bm m_j(u)$, including current, CV, CA, IMM,
robust IMM, or UKF, a valid bound is
\begin{equation}
 \norm{\bm x(t_b+u)-\bm m_j(u)}
 \leq r(u)+\norm{\bm m_j(u)-\bm c(u)}.
 \label{eq:recenter}
\end{equation}
\end{theorem}
\begin{proof}
Translate $t_b$ to zero. The error of extrapolating the two true positions
is the integral of acceleration against the nonnegative kernel
\begin{equation*}
 K(s)=\begin{cases}u(s+\Delta)/\Delta,&-\Delta\leq s\leq0,\\
 u-s,&0\leq s\leq u.\end{cases}
\end{equation*}
Its integral is $u(\Delta+u)/2$. The two measurement coefficients have
absolute values $1+u/\Delta$ and $u/\Delta$, giving
\eqref{eq:physical-radius} by the triangle inequality. In one dimension,
$x(t)=At^2/2$, $y_a=x(-\Delta)+\varepsilon_a$, and
$y_b=x(0)-\varepsilon_b$ attain equality. A final triangle inequality
gives \eqref{eq:recenter}.
\end{proof}

For equal error limits $\varepsilon$, the growth and noise sensitivity are
\begin{align}
 r(u)&=\varepsilon+b u+\tfrac12Au^2,
 &b&=2\varepsilon/\Delta+A\Delta/2,\\
 r'(u)&=b+Au,
 &\partial_\varepsilon r&=1+2u/\Delta.
 \label{eq:breakdown-rate}
\end{align}
For a tolerated radius $d>\varepsilon$ and $A>0$, this pair certificate
expires at the explicit horizon
\begin{equation}
 u_d=\frac{2(d-\varepsilon)}{b+\sqrt{b^2+2A(d-\varepsilon)}}.
 \label{eq:expiry}
\end{equation}
This is loss of the stated certificate, not a prediction that error
actually reaches $d$. For $A=0$ and $b>0$, $u_d=(d-\varepsilon)/b$.
For a different estimator center, the offset in \eqref{eq:recenter} must
also be included before asserting a deadline. The noise-amplification
and curvature terms balance at $\Delta=2\sqrt{\varepsilon/A}$ when both
parameters are positive. A shorter observation gap alone can therefore
worsen this two-point bound.

Every observation pair supplies a ball; their intersection retains all
such constraints. Appending observations while retaining the old balls
cannot enlarge that set. Sample count alone supplies no universal
$n^{-1/2}$ error law: a shared unknown sensor bias need not average away.
For a fixed history of $n$ isotropic Gaussian observations with known
marginal standard deviations, choosing
\begin{equation}
 \varepsilon_i=\sigma_i\sqrt{6\log(6n/\delta)}
\end{equation}
makes all observation limits hold with probability at least $1-\delta$,
by coordinate tails and a union bound. Independence is unnecessary for
this union bound, but the marginal error laws must be valid. On that
event, the physical balls hold simultaneously at all future times for
which the acceleration assumption holds. Adaptive stopping requires an
appropriate conditional or anytime observation guarantee; the fixed-$n$
argument alone does not justify it.

Even perfect knowledge of past position and velocity cannot remove the
worst-case radius $Au^2/2$: two continuations with opposite constant
accelerations share the entire past and separate by $Au^2$. Any common
prediction errs by at least half that separation on one continuation.
Without a future-motion restriction there is no finite uniform bound.
Consequently none of the compared estimators has a better worst-case
horizon order under this common assumption. Differences in their
recentered bounds remain possible, but have not been measured here.

The finite current model itself has velocity jumps at cell boundaries,
and the acquisition fixtures use linear interpolation on a fine grid.
Neither literally satisfies the bounded-acceleration premise in general.
We therefore do not attach an acceleration constant or certified physical
radius to those experiments. The physical theorem states what additional
motion and sensor information would be needed. Conversely, a smaller
model covariance, in this method or in a Kalman-family comparator, does
not establish a stronger physical guarantee.

\section{Results}
\label{sec:results}

\subsection{Comparison with the Earlier Tracking Controls}

\begin{table}[t]
\centering
\caption{Two-second endpoint RMSE (m) on the matched forecast screen.
C1: 65 observations, $\sigma=0.35$ m; C2: 65 observations,
$\sigma=1.4$ m; C3: 9 observations, $\sigma=0.35$ m.
Each column contains 30 paired cases.}
\label{tab:forecast}
\small
\begin{tabular}{lrrr}
\toprule Method & C1 & C2 & C3\\\midrule
\inputrows{tables/forecast_accuracy.tex}
\bottomrule
\end{tabular}
\end{table}

Table~\ref{tab:forecast} includes every method in the matched current
screen. At low noise with 65 observations, relational-current endpoint
RMSE is 2.958 m, compared with 3.750 for CV, 6.319 for CA, 3.443 for IMM,
3.269 for robust IMM, and 4.243 for the turn UKF. The paired current/robust
IMM ratio is 0.905 with interval $[0.837,0.981]$; the current/IMM ratio is
0.859 with interval $[0.779,0.952]$.

The same ordering does not hold in every condition. At high noise,
robust IMM reaches 4.806 m versus 4.891 m for the current mixture; their
ratio is 1.018 with interval $[0.921,1.124]$. In the sparse low-noise
condition, current RMSE is 3.788 m versus 3.940 m for robust IMM, with
ratio 0.961 and interval $[0.906,1.021]$. These intervals do not establish
a consistent advantage over robust IMM. Figure~\ref{fig:paired} shows
all five comparator ratios; the numerical intervals appear in
Appendix~\ref{app:intervals}.

\begin{figure*}[t]
\centering
\includegraphics[width=\textwidth]{paired_forecasts.pdf}
\caption{Current-model endpoint RMSE divided by comparator endpoint RMSE.
Values below one favor the current model. Bars are pointwise 95\%
paired percentile intervals from 4000 within-family resamples. The same
30 cases support each condition; the intervals do not cover arbitrary
motion families or simultaneous comparisons.}
\label{fig:paired}
\end{figure*}

The severed-current ablation increases C1 error to 4.868 m. A single
1.5 s length gives 3.556 m. The corresponding paired mixture/ablation
ratios are 0.608 $[0.537,0.675]$ and 0.832 $[0.757,0.908]$.
Severing establishes that the declared past--future covariance matters
for these forecasts. It does not establish uniqueness of that
continuation law. At high noise the single-length model reaches
4.762 m, slightly below the mixture's 4.891 m.

\begin{table}[t]
\centering
\caption{Endpoints contained in nominal 95\% moment ellipsoids.
These are empirical diagnostics, not exact mixture coverage guarantees.}
\label{tab:coverage}
\small
\begin{tabular}{lrrr}
\toprule Method & C1 & C2 & C3\\\midrule
\inputrows{tables/forecast_coverage.tex}
\bottomrule
\end{tabular}
\end{table}

Table~\ref{tab:coverage} reports containment using
$(x-\bar\mu)^\trans\Sigma^{-1}(x-\bar\mu)\leq7.814727903$.
The current mixture contains 30/30, 28/30, and 30/30 endpoints, with
mean equal-volume ellipsoid radii 5.77, 8.42, and 6.61 m. A Gaussian
mixture is not itself generally Gaussian, and these trajectories were
not drawn from the current prior. Consequently the chi-square threshold
does not certify 95\% posterior coverage here. Good endpoint RMSE and
large uncertainty regions must be assessed together.

\subsection{Measured Degradation Under Noise and Sparsity}

Table~\ref{tab:degradation} compares each method with its own C1 error.
The current mixture's RMSE rises 65.4\% when sensor standard deviation
quadruples, and 28.1\% when observations decrease from 65 to 9.
Robust IMM rises 47.0\% and 20.5\%, respectively. The current model thus
has larger relative degradation point estimates despite its lower C1
error and its still slightly lower C3 error. A lower starting error can
also increase a relative-growth ratio. These are two finite changes,
not fitted scaling exponents or formal breakdown laws; overlapping
pointwise intervals do not test the difference between methods' growth.
The earlier geometric prototype's degradation results concern a different
estimator and cannot establish a current-model rate.

\begin{table}[t]
\centering
\caption{Within-method endpoint RMSE inflation relative to C1, with
pointwise 95\% paired percentile intervals (4000 resamples within each
of five fixed families). Conditions are defined in Table~\ref{tab:forecast}.
An inflation of 1.654 means 65.4\% more RMSE.}
\label{tab:degradation}
\small\setlength{\tabcolsep}{3pt}
\begin{tabular}{lrr}
\toprule Method & C2/C1 [interval] & C3/C1 [interval]\\\midrule
\inputrows{tables/degradation.tex}
\bottomrule
\end{tabular}
\end{table}

\subsection{Acquisition Benefit and Settling}

\begin{table}[t]
\centering
\caption{Acquisition count and post-crossing accuracy.
Counts cover 12 s; RMSE uses the first six post-crossing seconds.
Rates are in Hz, errors in meters.}
\label{tab:acquisition}
\small\setlength{\tabcolsep}{3.5pt}
\begin{tabular}{rlrrr}
\toprule $\sigma$ & Rate profile & Samples & Position & 2 s forecast\\\midrule
\inputrows{tables/acquisition.tex}
\bottomrule
\end{tabular}
\end{table}

Added observations reduce current error and two-second prediction error
on these crossing fixtures (Table~\ref{tab:acquisition} and
Fig.~\ref{fig:acquisition}). At $\sigma=0.35$ m, adaptive 2-to-32 Hz
uses 167.8 samples on average, compared with 383.9 for constant 32 Hz.
Whole post-crossing position RMSE is 0.307 versus 0.209 m, so the
reduced acquisition count has a measurable accuracy cost over that
complete window.

\begin{table}[t]
\centering
\caption{Settling after the true crossing: position RMSE (m).
The true crossing defines evaluation windows only.}
\label{tab:settling}
\small
\begin{tabular}{rlrr}
\toprule $\sigma$ & Rate profile & First second & Seconds 3--6\\\midrule
\inputrows{tables/settling.tex}
\bottomrule
\end{tabular}
\end{table}

After three seconds, adaptive and always-dense tracking are close
(Table~\ref{tab:settling}). Their later RMSE ratio is 0.996 with interval
$[0.958,1.034]$ at low noise and 1.001 with interval $[0.951,1.053]$ at
high noise. The first-second ratios are 2.706 and 2.359, respectively.
Dense acquisition has already accumulated more information when the
boundary is crossed; adaptive acquisition pays for its sparse approach
with this settling interval.

Individual mixture corrections also become smaller as history grows.
At low noise, the later mean correction decreases from 0.702 m under
constant 2 Hz to 0.077 m under adaptive 2-to-32 Hz. The measured
projection of the correction onto the innovation decreases from 0.716
to 0.129. This projection describes the mixture correction; it does not impose
a clipping coefficient.

\begin{figure*}[t]
\centering
\includegraphics[width=\textwidth]{acquisition_accuracy.pdf}
\caption{Acquisition count versus post-crossing accuracy, averaged over
40 cases per noise level. Connected filled markers run from constant
2 Hz through the four adaptive target rates. Open markers are constant
32 Hz. All policies use the same physical evaluation grid. The same
accuracy values are retained by exact elimination; the figure shows
sampling effects rather than an optimization-induced accuracy change.}
\label{fig:acquisition}
\end{figure*}

\subsection{Exactness and Cost}

All 480 dense/compiled acquisition pairs have identical trigger times,
counts, timestamps, and observation arrays. Maximum absolute differences
are $7.29\times10^{-12}$ m for evaluated current positions,
$8.13\times10^{-12}$ m for two-second forecasts,
$6.26\times10^{-12}$ m for acquisition-event estimates, and
$1.37\times10^{-13}$ m for reported radius diagnostics.
These are numerical equivalence measurements, not improvements in
tracking accuracy or universal bitwise identities.

\begin{table}[t]
\centering
\caption{128-cell microbenchmark: median microseconds per call.
Forecast requests return current and two-second-ahead marginal moments.}
\label{tab:micro}
\small\setlength{\tabcolsep}{3.8pt}
\begin{tabular}{lrrr}
\toprule Implementation & Update & Position & Forecast\\\midrule
\inputrows{tables/micro.tex}
\bottomrule
\end{tabular}
\end{table}

The compiled update is 37.2 times faster than the original at 128 cells
(Table~\ref{tab:micro}). Moving the dense computation to real current
coordinates yields only 1.8 times, whereas the four-state NumPy
realization already yields 19.4 times. The state elimination is therefore
the main reduction. At 32, 128, and 512 cells, compiled updates take
11.99, 12.18, and 12.45 microseconds, versus 150.63, 452.68, and 5129.44
for dense Zak. The 128-cell cold initialization is 0.357 ms for the
compiled realization and 4.744 ms for dense Zak; with shared tables,
warm acquisition initialization is approximately 0.028 ms.

The small compiled CV cost control takes 5.39 microseconds per update
and 11.16 per two-query forecast. The five-component current model costs
2.26 and 2.02 times those operations, respectively. The NumPy CV control
takes 7.28 and 10.39 microseconds. Compilation therefore does not remove
all interface overhead, and neither comparison proves a lower bound
on achievable native latency. The CV control is independently checked
against its NumPy equations over 160 irregular updates and 320 queries.

\begin{table}[t]
\centering
\caption{Whole online compute over a 12 s trajectory, $\sigma=0.35$ m.
Means over 40 cases; costs include scheduling and readout, but not sensor
acquisition. Each pair follows identical acquisition decisions.}
\label{tab:online}
\small
\begin{tabular}{lrrr}
\toprule Rate profile & Dense (ms) & Compiled (ms) & Speedup\\\midrule
\inputrows{tables/online.tex}
\bottomrule
\end{tabular}
\end{table}

Whole-loop effects differ from isolated updates
(Table~\ref{tab:online} and Fig.~\ref{fig:cost}). Adaptive 2-to-32 Hz
falls from 114.47 to 12.55 ms, a 9.12-fold reduction. Assimilation falls
from 78.09 to 2.29 ms, while the unchanged scheduler takes 6.73 ms and
becomes the largest component. Constant 32 Hz falls from 218.18 to
9.58 ms. Thus the optimized adaptive loop costs \emph{more CPU} than
always-dense acquisition in this implementation, despite using only
43.7\% as many samples. This reversal is a consequence of inexpensive
filtering and nontrivial rate inversion. It must not be hidden by
reporting assimilation cost alone. The high-noise runs show the same
cost pattern.

\begin{figure*}[t]
\centering
\includegraphics[width=\textwidth]{exact_cost.pdf}
\caption{Left: update scaling with represented cell count, in double
precision. Right: whole online acquisition cost at $\sigma=0.35$ m;
each pair shows dense (left) and compiled (right). Scheduling is
unchanged and dominates the optimized adaptive loop. Microbenchmarks
use 31 repeats except 512-cell dense cases (five); whole-loop bars
average 40 matched cases per policy.}
\label{fig:cost}
\end{figure*}

\section{Interpretation and Limitations}

The computational and statistical claims have different scopes.
Theorem~\ref{thm:equivalence} establishes equality for the declared
finite Gaussian model. The synthetic comparisons ask whether that model
is useful outside its own prior distribution. Favorable low-noise
forecast errors cannot extend the theorem into a universal prediction
guarantee. Unknown future commands, correlated sensor drift, and noise
misspecification remain substantive limitations.

The measured current model is competitive with the selected conventional
implementations, but it does not decisively dominate robust IMM in every
condition. The earlier baseline forecasts also carry Monte Carlo error.
A more extensive benchmark would require external trajectories,
additional contemporary implementations, matched sensor assumptions,
and a common analytical or adequately converged forecast interface.
The separate simple CV timing control cannot supply those missing
accuracy comparisons.

The representation horizon is fixed, not periodically wrapped or
extended implicitly. Five 64-to-128-cell checks in the forecast screen
change endpoint means by less than 0.022 m; four 128-to-256-cell
acquisition checks change position RMSE by less than 0.00013 m and
forecast RMSE by less than 0.00084 m. These limited checks do not prove
a uniform continuum limit. Very fine cells also require care when
forming $Q_\ell$ by subtracting nearly equal matrices.

Twenty-three focused tests pass: 17 existing relational/acquisition
checks and six realization checks covering both NumPy and native
backends. They include kernel reconstruction, covariance stationarity,
8/32/128-cell streams, two coordinate systems, six sever configurations,
within-cell and irregular observations, evidence, fractional updates,
historical queries, full coefficient exports, and cache invalidation.
The archived dense oracle has the same source hash as the original
acquisition implementation. Five additional physical-certificate tests pass, including exact-rational
sharpness, a finite-turn path, retained constraints, recentering, and
invalid-input rejection. These test the conditional certificate separately.
Tests and replay establish the stated
implementation fidelity; they are not a formal floating-point proof.

\section{Conclusion}

A dense finite current history can be eliminated exactly when its
covariance admits a small Markov realization. For the specified
constant-plus-Mat\'ern-$3/2$ current model, accumulated position,
constant current, local current, and normalized current derivative
suffice. Shared spatial conditioning and evidence preserve the
five-component posterior while reducing a 128-cell update from
452.7 to 12.2 microseconds on the measured host.

The resulting computational scale is within a small factor of a simple
compiled CV filter, while the retained forecast comparison shows a
low-noise benefit over the measured robust IMM and no decisive benefit
in every other condition. Adaptive acquisition provides a separate
count--accuracy tradeoff with a settling interval. Once exact elimination
makes assimilation inexpensive, scheduling can determine the CPU cost
of that tradeoff. The representation, prediction law, numerical fidelity,
and acquisition policy must therefore be evaluated separately and then
accounted for together.

\appendices
\section{Complete Paired Forecast Intervals}
\label{app:intervals}

Table~\ref{tab:all-ci} gives the paired current/comparator RMSE ratios
under the protocol of Section~\ref{sec:forecast-protocol}. IMM and robust
IMM intervals retain the original reported bootstrap draws. CV, CA,
and UKF intervals are computed from the same retained case errors with
the same 4000-resample within-family procedure and a separately fixed
resampling seed. No method is retuned. All intervals are pointwise.

\begin{table}[h]
\centering
\caption{Current/comparator endpoint RMSE ratios and 95\% paired
percentile intervals. Condition definitions are in Table~\ref{tab:forecast}.}
\label{tab:all-ci}
\small
\begin{tabular}{llrr}
\toprule Condition & Comparator & Ratio & Interval\\\midrule
\inputrows{tables/paired_intervals.tex}
\bottomrule
\end{tabular}
\end{table}

\section{Earlier Geometric-Transport Confirmation}
\label{app:historical}

For completeness, Table~\ref{tab:historical} records the earlier
continuous-turn geometric-transport confirmation alongside the same
conventional filter families. This is a \emph{different cohort and
method}: 80 trajectories, eight seeds in each of five family/two
corruption strata, 6144 filtering particles for geometric transport,
and 256 forecast paths. It uses irregular observation intervals,
approximately two-second rolling forecasts, and corruption episodes
including drift, outliers, noise-scale changes, and missing samples.
The current model of Section~\ref{sec:model} was not evaluated in this
cohort, so its errors cannot be compared directly with this table.

The geometric prototype transports a particle distribution over speed,
direction, turn, position, and bias through finite rotation maps. Its
6144-particle confirmation obtained a 3.6\% endpoint RMSE reduction
against robust IMM, with a paired interval of $[-2.1\%,8.5\%]$.
That result did not establish an advantage over the strongest robust
control. Its crossing Brier score was also worse. These findings
motivate keeping accuracy, uncertainty, and boundary decisions distinct;
they are not additional evidence of current-model superiority.

\begin{table}[h]
\centering
\caption{Separate geometric confirmation cohort. Tracking/forecast
RMSE and energy score are in meters. Brier score concerns sampled-path
crossing of $x=0$. The current model is absent from this cohort.}
\label{tab:historical}
\small\setlength{\tabcolsep}{3pt}
\begin{tabular}{lrrrr}
\toprule Method & Track & $\sim$2 s & Energy & Brier\\\midrule
\inputrows{tables/historical.tex}
\bottomrule
\end{tabular}
\end{table}

\section{Reproduction and Artifact Map}
\label{app:artifact}
\begingroup\raggedright

The accompanying evidence archive retains the forecast, acquisition,
optimization, and geometric-confirmation records separately, with a
SHA-256 source manifest. The manuscript's tables and vector figures are
generated from these records by \texttt{generate\_assets.py}; rendering
performs no fitting or new estimator selection. The TeX bundle includes
all generated table inputs and figures, and uses standard packages
without shell escape or machine-specific fonts. It compiles without
access to the experiment repository.

The repository experiment directory is
\texttt{experiments/civilian\_transport}. Its
\texttt{relational\_reference.py} preserves the dense oracle;
\texttt{relational\_markov.py} specifies the four-state algebra;
\texttt{relational\_native.py} and \texttt{native/markov.cpp} expose
and implement the compiled operations. The runner
\texttt{run\_optimization\_m4.sh} builds the optional kernels, runs the
23 focused tests, performs the 480-case paired replay and microbenchmark,
and copies the result files back from the compute host. The standalone
cost control is \texttt{cost\_controls.py}.

The earlier forecast screen is reproduced by
\texttt{run\_relational\_m4.sh}; acquisition accuracy by
\texttt{run\_acquisition\_m4.sh}; and the separate geometric
confirmation by \texttt{run\_confirmation\_m4.sh}. Re-running through
the promoted streaming class changes computation cost but not the
specified finite prior. Original timing evidence remains archived
rather than being silently replaced by a rerun.

\par\endgroup

\begin{table}[h]
\centering
\caption{128-cell update variability: median and 10th--90th percentiles
of repeat-average microseconds per observation.}
\label{tab:variability}
\small
\begin{tabular}{lrr}
\toprule Implementation & Median & 10th--90th percentile\\\midrule
\inputrows{tables/micro_variability.tex}
\bottomrule
\end{tabular}
\end{table}

\begin{thebibliography}{99}
\bibitem{kalman1960}
R. E. Kalman, ``A new approach to linear filtering and prediction
problems,'' \emph{J. Basic Eng.}, vol. 82, no. 1, pp. 35--45, 1960,
doi: \href{https://doi.org/10.1115/1.3662552}{10.1115/1.3662552}.
\bibitem{hartikainen2010}
J. Hartikainen and S. S\"arkk\"a, ``Kalman filtering and smoothing
solutions to temporal Gaussian process regression models,'' in
\emph{Proc. IEEE Int. Workshop Mach. Learn. Signal Process.}, 2010,
pp. 379--384, doi: \href{https://doi.org/10.1109/MLSP.2010.5589113}{10.1109/MLSP.2010.5589113}.
\bibitem{corenflos2021}
A. Corenflos, Z. Zhao, and S. S\"arkk\"a, ``Temporal Gaussian process
regression in logarithmic time,'' arXiv:2102.09964, 2021.
[Online]. Available: \url{https://arxiv.org/abs/2102.09964}
\bibitem{blom1988}
H. A. P. Blom and Y. Bar-Shalom, ``The interacting multiple model
algorithm for systems with Markovian switching coefficients,''
\emph{IEEE Trans. Autom. Control}, vol. 33, no. 8, pp. 780--783, 1988,
doi: \href{https://doi.org/10.1109/9.1299}{10.1109/9.1299}.
\bibitem{julier2004}
S. J. Julier and J. K. Uhlmann, ``Unscented filtering and nonlinear
estimation,'' \emph{Proc. IEEE}, vol. 92, no. 3, pp. 401--422, 2004,
doi: \href{https://doi.org/10.1109/JPROC.2003.823141}{10.1109/JPROC.2003.823141}.
\bibitem{lizzio2026}
F. F. Lizzio, E. I. Trombetta, E. Capello, and Y. Fujisaki,
``Comparative evaluation of multiple-model Kalman filters for highly
maneuvering UAV tracking,'' \emph{Appl. Sci.}, vol. 16, no. 5,
Art. no. 2377, 2026, doi: \href{https://doi.org/10.3390/app16052377}{10.3390/app16052377}.
\bibitem{koller2021}
T. L. Koller and U. Frese, ``The interacting multiple model filter and
smoother on boxplus-manifolds,'' \emph{Sensors}, vol. 21, no. 12,
Art. no. 4164, 2021, doi: \href{https://doi.org/10.3390/s21124164}{10.3390/s21124164}.
\bibitem{zhang2017}
C. Zhang, A. Taghvaei, and P. G. Mehta, ``Feedback particle filter on
matrix Lie groups,'' arXiv:1701.02416, 2017.
[Online]. Available: \url{https://arxiv.org/abs/1701.02416}
\end{thebibliography}
\end{document}
